\chapter{Zadatak 1d --- rješenje u MATLAB-u}
\label{ch:z1-matlab}

\section{Što se ovdje traži}

Riješiti isti problem rješavačem i \emph{usporediti} dobivene primarne i dualne
varijable s analitičkim rezultatom iz (b). Naglasak je na usporedbi --- sam
poziv funkcije nije poanta.

\begin{zadatak}{Zadatak 1, podzadatak d)}
d) Riješite problem u MATLABu (quadprog i/ili yalmip) --- odredite primarne i
dualne varijable u optimumu te ih usporedite s rezultatima (Lagrangeovim
množiteljima) iz (b).
\end{zadatak}

\begin{predznanje}
\begin{itemize}[nosep,leftmargin=*]
  \item Standardni oblik QP-a $\min \tfrac12 x\trans Qx + c\trans x$ uz $Gx \leq h$ --- \S\ref{sec:z1-standardni-oblik}.
  \item Matrice $Q$, $c$, $G$, $h$ --- \eqref{eq:z1-Qc} i \eqref{eq:z1-Gh}.
  \item Analitičko rješenje s kojim se uspoređuje --- \eqref{eq:z1b-rjesenje}.
\end{itemize}
\end{predznanje}

\section{Zašto \texttt{quadprog} traži baš te argumente}
\label{sec:z1d-sucelje}

Sučelje funkcije je
\begin{center}
\texttt{quadprog(H, f, A, b, Aeq, beq, lb, ub, x0, options)}
\end{center}
i minimizira $\tfrac12 x\trans H x + f\trans x$ uz $Ax \leq b$. To je \emph{točno}
standardni oblik s Predavanja 5, pa se matrice predaju izravno, bez ikakve
pretvorbe:
\begin{equation}
\texttt{H} = Q, \quad \texttt{f} = c, \quad \texttt{A} = G, \quad \texttt{b} = h .
\end{equation}

\begin{zamka}[Zamka: MATLAB-ov $\tfrac12$ i ,,zaboravljeni`` faktor]
\texttt{quadprog} \emph{sam} podrazumijeva faktor $\tfrac12$ ispred kvadratnog
člana. Ako se preda matrica bez tog faktora --- dakle
$\begin{bsmallmatrix}1 & -0{,}15\\ -0{,}15 & 2\end{bsmallmatrix}$ umjesto $Q$ ---
rješavač će riješiti dvostruko manji kvadratni član i vratiti krivi optimum. Ovo
je najčešća pogreška pri prvom pozivu, a ne javlja se kao greška nego kao
naizgled uredan, ali pogrešan rezultat. Kontrola: usporedi
\texttt{fval} s $f(\xopt)$ izračunatim ručno.
\end{zamka}

\begin{zamka}[Zamka: $h$, a ne $-h$]
Sučelje traži $Ax \leq b$, dakle desnu stranu \emph{prije} prebacivanja. Naš
$h = (-10,\; 3)$ ide izravno. Ako se preda $(10,\, -3)$, dopustivi skup je
zrcaljen i rješenje je besmisleno.
\end{zamka}

\section{Kako je to izgledalo na vježbama}
\label{sec:z1d-vjezbe}

\begin{izvor}{Vježbe 3 --- Primjer 2, slajd 25/36 (PDF str.~30)}
\slajd[0.86]{v3-primjer2-quadprog}

\textbf{Što je na slajdu.} Snimka MATLAB-ove sesije: poziv \texttt{quadprog} za
profesorov primjer, s ispisanim rješenjem i množiteljima. Kod je strukturno
identičan našemu --- definiraju se $Q$, $c$, $G$, $h$, pa se pozove rješavač.

\textbf{Što iz toga treba zapamtiti:} profesor \emph{pokazuje kod na slajdu}.
Znači da se na obrani očekuje da znaš objasniti svaki argument poziva, a ne samo
rezultat.
\end{izvor}

\begin{izvor}{Vježbe 3 --- Primjer 2, slajd 28/36 (PDF str.~33)}
\slajd[0.86]{v3-primjer2-yalmip}

\textbf{Što je na slajdu.} Isti problem riješen preko YALMIP-a, koji kao rješavač
koristi \texttt{quadprog}. Varijable se deklariraju sa \texttt{sdpvar},
ograničenja se skupljaju u listu, pa se pozove \texttt{optimize}.

\textbf{Zašto to zadatak spominje kao alternativu:} YALMIP je modelirajući sloj,
ne rješavač. Za problem koji je već u standardnom obliku ne donosi ništa; postaje
nezamjenjiv u Zadatku 4, gdje ograničenja nisu $Gx \leq h$ nego matrične
nejednakosti.
\end{izvor}

\section{Kod i stvarni ispis}
\label{sec:z1d-kod}

\begin{seminar}{Seminar v1.2, §2.5 --- ,,(d) Numeričko rješenje u MATLAB-u``}
Problem se u MATLAB-u rješava funkcijom \texttt{quadprog}, koja minimizira
$\tfrac{1}{2}x^{\mathsf{T}}Qx + c^{\mathsf{T}}x$ uz $Gx \leq h$, pa se matrice
predaju izravno.

\begin{lstlisting}[style=matlabstyle]
Q = [2.0, -0.3; -0.3, 4.0];
c = [2.0; -3.0];
G = [-1.0, -1.0; 1.0, -1.0];
h = [-10.0; 3.0];

opts = optimoptions('quadprog', 'Display', 'none', ...
                    'Algorithm', 'interior-point-convex');
[x_qp, f_qp, ef_qp, ~, lam_qp] = quadprog(Q, c, G, h, [], [], [], [], [], opts);
\end{lstlisting}

\begin{lstlisting}[style=outputstyle]
exitflag        = 1
x_opt           = [5.757576; 4.242424]
p*  = fval      = 60.606061
lambda.ineqlin  = [12.242424; 0.000000]
\end{lstlisting}

Primarne i dualne varijable dobivene rješavačem podudaraju se s analitičkim
rezultatom do na tolerancije rješavača. Polje \texttt{lambda.ineqlin} sadrži
upravo Lagrangeove množitelje $\mu_1^*$ i $\mu_2^*$ iz KKT uvjeta.
\end{seminar}

\subsection{Što koji redak znači}

\begin{longtable}{@{}p{0.30\linewidth} p{0.66\linewidth}@{}}
\toprule
\textbf{Redak} & \textbf{Značenje} \\
\midrule
\endhead
\texttt{'Display','none'} & Isključuje ispis rješavača, jer se sve što treba
ispisuje ručno. \\[5pt]
\texttt{'Algorithm','interior-\allowbreak point-convex'} & Metoda unutarnje točke
za konveksne QP-ove. Ne nabraja aktivne skupove kao mi u (b), nego se optimumu
približava iznutra i aktivni skup ,,otkriva`` u limesu. Zato su
$\mu_2^{\star}$ i $g_2$ jednaki nuli tek do na toleranciju. \\[5pt]
\texttt{ef\_qp} (\texttt{exitflag}) & $1$ znači da je rješavač konvergirao.
Svaka druga vrijednost čini rezultat neupotrebljivim i mora se provjeriti prije
nego se brojke uopće pogledaju. \\[5pt]
\texttt{lam\_qp.ineqlin} & Vektor Lagrangeovih množitelja ograničenja
nejednakosti --- upravo $\mu$ iz KKT sustava. Postoje i polja
\texttt{eqlin}, \texttt{lower}, \texttt{upper} za ostale vrste ograničenja; ovdje
su prazna. \\[5pt]
\texttt{\textasciitilde} kao četvrti izlaz & Preskočeni izlaz
\texttt{output} (statistika iteracija), koji nam ne treba. \\
\bottomrule
\end{longtable}

\section{Usporedba, koja je zapravo zadatak}
\label{sec:z1d-usporedba}

\begin{seminar}{Seminar v1.2, Tablica 2.2}
\begin{center}
\begin{tabular}{lccc}
\toprule
\textbf{Veličina} & \textbf{Analitički} & \texttt{quadprog} & \textbf{Odstupanje} \\
\midrule
$x_1^*$   & $190/33 = 5.757576$  & $5.757576$  & $< 10^{-10}$ \\
$x_2^*$   & $140/33 = 4.242424$  & $4.242424$  & $< 10^{-10}$ \\
$\mu_1^*$ & $404/33 = 12.242424$ & $12.242424$ & $< 10^{-10}$ \\
$\mu_2^*$ & $0$                  & $0$         & --- \\
$p^*$     & $2000/33 = 60.606061$& $60.606061$ & $< 10^{-10}$ \\
\bottomrule
\end{tabular}
\end{center}
\end{seminar}

Iz punog ispisa programa vidi se i norma razlike:

\begin{lstlisting}[style=outputstyle]
|x_qp - x_analiticki| = 7.155e-11
\end{lstlisting}

\begin{korak}[Što se ovom usporedbom zapravo pokazalo]
\begin{enumerate}[leftmargin=*]
  \item Analitičko rješavanje KKT sustava i numerički rješavač daju istu točku ---
        dakle nijedan od dva puta nema grešku.
  \item Množitelji koje vraća rješavač \emph{jesu} KKT množitelji. To nije
        samorazumljivo: rješavač ih ne računa iz naše tablice slučajeva, nego kao
        dualne varijable svoje metode unutarnje točke. Da se poklapaju, potvrda
        je da su oba pristupa opis istog matematičkog objekta.
  \item Odstupanje reda $10^{-11}$ je tolerancija rješavača, a ne fizikalna
        veličina. Zato se u seminaru rezultat navodi kao razlomak, s decimalama
        u zagradi.
\end{enumerate}
\end{korak}

\section{Modelski odgovori}
\label{sec:z1d-odgovori}

\begin{pitanje}{Zašto \texttt{quadprog} vraća $\mu_2 = 0$, a ne neki mali broj poput $10^{-12}$?}
\textbf{Odgovor.} Vraća ga do na toleranciju; ispisan je kao \texttt{0.000000}
jer se ispisuje na šest decimala. Metoda unutarnje točke u limesu tjera umnožak
$\mu_j g_j$ prema nuli, pa za neaktivno ograničenje $\mu_j$ pada prema nuli
zajedno s parametrom barijere. Da je algoritam bio tipa aktivnog skupa,
$\mu_2$ bi bio točno nula po konstrukciji.
\end{pitanje}

\begin{pitanje}{Biste li isti rezultat dobili preko YALMIP-a?}
\textbf{Odgovor.} Da. YALMIP je modelirajući sloj --- problem se zapiše
deklarativno, a on ga proslijedi rješavaču, među ostalim i \texttt{quadprog}-u.
Za ovako mali QP nema razlike u rezultatu; YALMIP se isplati kad ograničenja nisu
u standardnom obliku, kao u Zadatku 4, gdje se rješavaju linearne matrične
nejednakosti.
\end{pitanje}

\begin{pitanje}{Što ako \texttt{exitflag} nije 1?}
\textbf{Odgovor.} Tada se rezultat ne smije koristiti. Vrijednost $-2$ znači da
problem nije dopustiv, $-3$ da je neograničen, $0$ da je premašen broj
iteracija. Ovdje unaprijed znamo da je problem dopustiv (Slaterova točka) i
ograničen odozdo (koercivnost), pa bi svaka od tih vrijednosti značila grešku u
predaji matrica.
\end{pitanje}

\begin{izgovor}
,,U MATLAB-u sam koristio \texttt{quadprog}, jer njegovo sučelje već jest
standardni oblik s predavanja --- minimizira jedna polovina iks transponirano ha
iks plus ef transponirano iks uz a iks manje ili jednako be --- pa sam matrice
predao izravno, bez pretvorbe. Bitno je da faktor jedna polovina rješavač
podrazumijeva, pa se predaje puna matrica ku.

Rješavač je vratio iks pet zarez sedam pet sedam pet sedam šest i četiri zarez
dva četiri dva četiri, optimalnu vrijednost šezdeset zarez šest nula šest, i
množitelje dvanaest zarez dva četiri i nulu. To se s analitičkim rješenjem iz
podzadatka b) poklapa do na jedanaestu decimalu.

Time je pokazano i nešto više od podudarnosti brojeva: polje
\texttt{lambda.ineqlin} koje rješavač vraća sadrži upravo Lagrangeove
množitelje iz KKT uvjeta, iako ih algoritam dobiva sasvim drugim putem --- kao
dualne varijable metode unutarnje točke.``
\end{izgovor}
